\subsubsection{\texttt{spmak}}

\texttt{spmak}函数根据用户给出的信息创建一个 B 格式的样条曲线。，其用法如下:

\texttt{sp = spmak(knots,coefs)},该函数返回值为一个结构体(B 格式)，其中

\begin{itemize}
	\item \texttt{knots} 为 B 样条节点,这里要求其为单增非降的一维向量；
	\item \texttt{coefs} 为 B 样条的系数,要求其为一维向量，且长度必须小于 \texttt{knots} 的长度；
	\item 样条曲线的阶数 \texttt{k} 满足 \texttt{k = length(knots) - length(coefs)}
	\begin{itemize}
		\item 需要注意，节点的重数需要满足\texttt{$\leq k$}。如果对应的B-spline只有一个不同的节点，即\texttt{knots(j) = knots(j+k)}，那么系数\texttt{coefs(:,j)}将被直接忽略。
	\end{itemize}
	\item \texttt{sp} 为一结构体，包含如下字段: \texttt{form} (样条的形式，在这里为 \texttt{'B-'} 格式)、\texttt{knots} (样条节点的位置，为一维向量)、\texttt{coefs}(B样条基函数系数构成的一维数向量)、\texttt{number}(B样条基函数的个数)、\texttt{order} (样条函数的阶数)、\texttt{dim} (目标函数的维数);
\end{itemize}

如给定 B 样条节点 $\{0,0,1,2,3,4,5,5\}$, 可以得到三个4次的 B 样条，若这三个 B 样条的系数分别为1，4和-2，则我们可以得到对应的
B 格式样条曲线。
\par 输入为：
\begin{lstlisting}[language = matlab]
        knots = [0 0 1 2 3 4 5 5]; 
        coefs = [1 4 -2]; 
        sp = spmak(knots,coefs)
\end{lstlisting}

输出为：
\begin{lstlisting}[language = matlab ]
     coefs : [1,4,-2]
       dim : 1
      form : 'B-'
     knots : [0,0,1,2,3,4,5,5]
    number : 3
     order : 5
\end{lstlisting}
